\documentclass[11pt,a4paper]{article}
\usepackage[margin=0.85in,top=1in,bottom=1in]{geometry}
\usepackage{amsmath,amssymb,amsfonts}
\usepackage{graphicx}
\usepackage{float}
\usepackage{enumitem}
\usepackage{microtype}
\usepackage{caption}
\usepackage[hidelinks]{hyperref}
\usepackage[most,breakable]{tcolorbox}
\tcbuselibrary{breakable,skins}
\usepackage[utf8]{inputenc}
\usepackage[T1]{fontenc}
\usepackage{lmodern}
\captionsetup{font=small,labelfont=bf,skip=4pt}
\setlist{leftmargin=1.5em,topsep=2pt}

%% Breakable card — prevents content cutoff across pages
\newtcolorbox{solcard}[3]{
  breakable, enhanced,
  colback=white, colframe=gray!50, arc=2pt,
  title={\textbf{#1}\hfill\textsf{\small #3\,pts}},
  fonttitle=\normalsize\sffamily\bfseries,
  before upper={\small\textit{#2}\par\smallskip\hrule height 0.4pt\smallskip},
  top=4pt, bottom=6pt, left=7pt, right=7pt,
  parbox=false,
}

\title{\textbf{Gyroscopic forces and parasitic friction --- Solution}}
\author{\normalsize Y23 (2023) — English translation}
\date{}
\begin{document}\maketitle\vspace{-0.5em}
\begin{solcard}{A1}{Let us consider a ring charged uniformly around its perimeter with a charge $Q$ and rotating around its axis with an angular velocity $\Omega$. Find the equivalent current $i$ in the ring.}{0.30}

First solution: During the rotation time $T=2\pi/\Omega$, each elementary charge passes through a point in the ring, from where: $$i=\cfrac{Q}{T} $$


Second solution: The result can also be obtained using the expression: $$i=\lambda v $$ where $v=\Omega R$ - speed of ring elements. Both expressions lead to the answer:

\par\smallskip\noindent\textbf{Answer:} $$i=\cfrac{\Omega Q}{2\pi} $$
\end{solcard}

\begin{solcard}{A2}{Determine the magnetic moment of the ball $\vec{m}$. Express your answer in terms of $\vec{\omega}$, $q$ and $r$.}{1.00}

First solution: The magnetic moment of the system is determined by the relation: $$\vec{m}=\int \cfrac{\bigl[\vec{r}\times\vec{v}\bigr]dq}{2} $$ where $dq$ is an elementary charge, and $\vec{r}$ and $\vec{v}$ - it radius-vector and velocity correspondingly. Note, that angular momentum mechanical system be determined relation: $$\vec{L}=\int \bigl[\vec{r}\times\vec{v}\bigr] dM $$ where $dM $ - elementary mass, and $\vec{r}$ and $\vec{v}$ are its radius-vector and velocity correspondingly. If the distributions of charge and mass are similar, we arrive at the following relation: $$\vec{m}=\cfrac{\vec{L}q}{2M}=\cfrac{Iq\vec{\omega}}{2M} $$ where $I=2Mr^2/5$ - moment of inertia of a homogeneous ball relative to the diameter.


$$\vec{m}=\int \cfrac{\bigl[\vec{r}\times\vec{v}\bigr]dq}{2} $$


where $dq$ is the elementary charge, and $\vec{r}$ and $\vec{v}$ are its radius vector and speed, respectively.


Note that the angular momentum of a mechanical system is determined by the relation:


$$\vec{L}=\int \bigl[\vec{r}\times\vec{v}\bigr] dM $$


where $dM$ is the elementary mass, and $\vec{r}$ and $\vec{v}$ are its radius vector and speed, respectively.


If the charge and mass distributions are similar, then we arrive at the following relationship:


$$\vec{m}=\cfrac{\vec{L}q}{2M}=\cfrac{Iq\vec{\omega}}{2M} $$


where $I=2Mr^2/5$ is the moment of inertia of a homogeneous ball relative to the diameter.


Second solution: consider a disk with radius $R$ and thickness $h$ with charge density $\rho$ and rotating with angular velocity $\vec{\omega}$ around its rotation axis. Select a ring with radius $r$ and thickness $dr$. Then for the magnetic moment of the ring we obtain: $$d\vec{m}_\text{further}=\vec{n}S\cdot\cfrac{\omega dQ}{2\pi} $$ where $S=\pi{r}^2$ - area ring. Transform: $$d\vec{m}_{\text{further}}=\vec{\omega}\cdot\cfrac{2\pi{r}^3drh\rho}{2}\Rightarrow{\vec{m}_\text{further}=\cfrac{\pi{R}^4h\rho\vec{\omega}}{2}} $$ Transition to sphere. Be characterizing $h$ and $R$ angle $\theta$ with axis of symmetry sphere. Then: $$R=r\sin\theta\qquad h=-r\cos\theta $$ Then for element magnetic moment sphere have: $$d\vec{m}=\cfrac{\pi{r}^5\rho\vec{\omega}}{4}\sin^5\theta d\theta\Rightarrow{\vec{m}=\cfrac{\pi{r}^5\rho\vec{\omega}}{4}\int\limits_0^{\pi}\sin^5\theta d\theta}=\cfrac{\pi{r}^5\rho\vec{\omega}}{4}\int\limits_{-1}^{1}(1-2t^2+t^4)dt=\cfrac{4\pi{r}^5\rho\vec{\omega}}{15} $$ For answer on question remain substitute $\rho=3q/4\pi{r}^3$.


$$d\vec{m}_\text{further}=\vec{n}S\cdot\cfrac{\omega dQ}{2\pi} $$


where $S=\pi{r}^2$ is the area of the ring. Let's convert:


$$d\vec{m}_{\text{further}}=\vec{\omega}\cdot\cfrac{2\pi{r}^3drh\rho}{2}\Rightarrow{\vec{m}_\text{further}=\cfrac{\pi{R}^4h\rho\vec{\omega}}{2}} $$


Let's move on to the ball. We will characterize $h$ and $R$ by the angle $\theta$ with the symmetry axis of the ball. Then:


$$R=r\sin\theta\qquad h=-r\cos\theta $$


Then for the element of the magnetic moment of the ball we have:


$$d\vec{m}=\cfrac{\pi{r}^5\rho\vec{\omega}}{4}\sin^5\theta d\theta\Rightarrow{\vec{m}=\cfrac{\pi{r}^5\rho\vec{\omega}}{4}\int\limits_0^{\pi}\sin^5\theta d\theta}=\cfrac{\pi{r}^5\rho\vec{\omega}}{4}\int\limits_{-1}^{1}(1-2t^2+t^4)dt=\cfrac{4\pi{r}^5\rho\vec{\omega}}{15} $$


To answer the question, all that remains is to substitute $\rho=3q/4\pi{r}^3$.


Third solution: consider a thin-walled spherical shell with radius $R$ and thickness $dR$, rotating with angular velocity $\vec{\omega}$. Let us select a ring limited by the angles $\theta$ and $\theta+d\theta$ of the spherical coordinate system. The magnetic moment of the selected ring is equal to: $$d\vec{m}_\text{sph}=\cfrac{dQ\vec{\omega}}{2\pi}\cdot \pi{R}^2\sin^2\theta=\vec{\omega}\cdot\cfrac{2\pi\rho R^2dR\sin\theta d\theta}{2\pi}\cdot\pi{R}^2\sin^2\theta=\pi\rho{R}^4dR\vec{\omega}\cdot\sin^3\theta d\theta $$ Integgives the expression: $$\vec{m}_\text{sph}=\pi\rho{R}^4dR\vec{\omega}\int\limits_0^{\pi}\sin^3\theta d\theta=\pi\rho{R}^4dR\vec{\omega}\int\limits_{-1}^{1}(1-t^2)dt=\cfrac{4\pi\rho R^4dR\vec{\omega}}{3} $$ Magnetic moment sphere be obtained integration magnetic moment by sphere: $$\vec{m}=\cfrac{4\pi\rho\vec{\omega}}{3}\int\limits_{0}^rR^4dR=\cfrac{4\pi\rho{r}^5\vec{\omega}}{15} $$ For answer on question remain substitute $\rho=3q/4\pi{r}^3$. Finally:


$$d\vec{m}_\text{sph}=\cfrac{dQ\vec{\omega}}{2\pi}\cdot \pi{R}^2\sin^2\theta=\vec{\omega}\cdot\cfrac{2\pi\rho R^2dR\sin\theta d\theta}{2\pi}\cdot\pi{R}^2\sin^2\theta=\pi\rho{R}^4dR\vec{\omega}\cdot\sin^3\theta d\theta $$


Let's integrate the resulting expression:


$$\vec{m}_\text{sph}=\pi\rho{R}^4dR\vec{\omega}\int\limits_0^{\pi}\sin^3\theta d\theta=\pi\rho{R}^4dR\vec{\omega}\int\limits_{-1}^{1}(1-t^2)dt=\cfrac{4\pi\rho R^4dR\vec{\omega}}{3} $$


The magnetic moment of a ball is obtained by integrating the magnetic moment over the spheres:


$$\vec{m}=\cfrac{4\pi\rho\vec{\omega}}{3}\int\limits_{0}^rR^4dR=\cfrac{4\pi\rho{r}^5\vec{\omega}}{15} $$


To answer the question, all that remains is to substitute $\rho=3q/4\pi{r}^3$.


Finally:

\par\smallskip\noindent\textbf{Answer:} $$\vec{m}=\cfrac{qr^2\vec{\omega}}{5} $$
\end{solcard}

\begin{solcard}{A3}{Find the moment of forces $\vec{M}$ acting on the ball relative to its center. Express your answer in terms of $\vec{\omega}$, $\vec{B}$, $q$ and $r$.}{0.20}

The moment of forces $\vec{M}$ acting on a coil with current in a uniform magnetic field is equal to: $$\vec{M}=\bigl[\vec{m}\times\vec{B}\bigr] $$from where:

\par\smallskip\noindent\textbf{Answer:} $$\vec{M}=\cfrac{qr^2\bigl[\vec{\omega}\times\vec{B}\bigr]}{5} $$
\end{solcard}

\begin{solcard}{B1}{Write down the condition for the ball not to slip on the plane in terms of $\vec{v}_C$, $\vec{\omega}$ and $\vec{r}$. Using the resulting expression, express the angular velocity components $\omega_x$ and $\omega_y$ in terms of $v_{Cx}$, $v_{Cy}$ and $r$. Differentiating the no-slip condition with respect to time, express $\vec{a}_C$ in terms of $\dot{\vec{\omega}}$ and $\vec{r}$.}{0.50}

If a rigid body rotates with angular velocity $\vec{\omega}$, and the speed of point $A$ is equal to $\vec{v}_A$, then the speed of point $B$ is given by the expression: $$\vec{v}_B=\vec{v}_A+\bigl[\vec{\omega}\times\overrightarrow{AB}\bigr] $$From here we obtain the condition of no slipping:


Let's project the resulting equation onto the $x$ and $y$ axes: $$v_{Cx}-\omega_yr=0\qquad v_{Cy}+\omega_xr=0 $$From here:


The vector $\vec{r}$ remains constant as the ball moves along the plane. Then differentiating, we get:

\par\smallskip\noindent\textbf{Answer:} $$\vec{v}_C+\bigl[\vec{\omega}\times\vec{r}\bigr]=0 $$
\end{solcard}

\begin{solcard}{B2}{Using the equation for the dynamics of rotational motion about the center of mass, express the angular acceleration of the ball $\dot{\vec{\omega}}$ in terms of $q$, $m$, $r$, $\vec{\omega}$, $\vec{r}$, $\vec{B}$ and $\vec{F}$.}{0.20}

Let us write down the equation for the dynamics of rotational motion relative to the center of mass: $$\cfrac{d\vec{L}_C}{dt}=\vec{M}_C $$For moment force relatively of the center of mass have: $$\vec{M}_C=\bigl[\vec{r}\times\vec{F}\bigr]+\cfrac{qR^2\bigl[\vec{\omega}\times\vec{B}\bigr]}{5} $$and for derivative moment momentum by time have: $$\cfrac{d\vec{L}_C}{dt}=I\dot{\vec{\omega}} $$from:

\par\smallskip\noindent\textbf{Answer:} $$\dot{\vec{\omega}}=\cfrac{5\bigl[\vec{r}\times\vec{F}\bigr]}{2mr^2}+\cfrac{q\bigl[\vec{\omega}\times\vec{B}\bigr]}{2m} $$
\end{solcard}

\begin{solcard}{B3}{Express the friction force $\vec{F}$ in terms of $m$, $q$, $\vec{B}$, $\vec{\omega}$, $\vec{r}$ and $\vec{a}_C$. Note: Use the double cross product property: $$\bigl[\vec{a}\times\bigl[\vec{b}\times\vec{c}\bigr]\bigr]=\vec{b}\bigl({~}\vec{a}\cdot\vec{c}{~}\bigr)-\vec{c}\bigl({~}\vec{a}\cdot\vec{b}{~}\bigr)$$}{0.20}

Note: Take advantage of the double cross product property:


$$\bigl[\vec{a}\times\bigl[\vec{b}\times\vec{c}\bigr]\bigr]=\vec{b}\bigl({~}\vec{a}\cdot\vec{c}{~}\bigr)-\vec{c}\bigl({~}\vec{a}\cdot\vec{b}{~}\bigr)$$


Consider the vector product with $\vec{r}$ terms in the equation obtained in step $\mathrm{B3}$: from:

\par\smallskip\noindent\textbf{Answer:} $$\vec{F}=-\cfrac{2m\vec{a}_C}{5}-\cfrac{q\vec{B}\bigl({~}\vec{r}\cdot\vec{\omega}{~}\bigr)}{5} $$
\end{solcard}

\begin{solcard}{B4}{Find the angular velocity component of the ball $\omega_z$ at the point with coordinate $x_C$ of its center. Express your answer in terms of $q$, $m$, $B$, $r$ and $x_C$.}{0.80}

Since the vectors $\vec{r}$ and $\vec{e}_z$ are collinear, the equation for the dynamics of rotational motion in projection onto the $z$ axis is written as follows: $$I\dot{\omega}_z=\cfrac{qr^2\bigl[\vec{\omega}\times\vec{B}\bigr]_z}{5}=-\cfrac{qr^2B\omega_y}{5}\Rightarrow{\dot{\omega}_z=-\cfrac{qB\omega_y}{2m}} $$On the other hand, $\omega_yr=v_x$, from which: $$\dot{\omega}_z=-\cfrac{qBv_x}{2mr} $$and finally:

\par\smallskip\noindent\textbf{Answer:} $$\omega_z=-\cfrac{qBx_C}{2mr} $$
\end{solcard}

\begin{solcard}{B5}{Prove that the trajectory of the center of the ball is rectilinear.}{0.40}

Let us write the theorem on the motion of the center of mass: $$m\vec{a}_C=m\vec{g}+\vec{N}+\vec{F}+q\bigl[\vec{v}_C\times\vec{B}\bigr] $$Substitute expression for $\vec{F}$, obtain in part $\mathrm{B3}$, obtain: $$m\vec{a}_C=m\vec{g}+\vec{N}+q\bigl[\vec{v}_C\times\vec{B}\bigr]-\cfrac{2m\vec{a}_C}{5}-\cfrac{q\vec{B}\bigl({~}\vec{r}\cdot\vec{\omega}{~}\bigr)}{5} $$from which: $$\vec{a}_C=\cfrac{5\vec{g}}{7}+\cfrac{5\vec{N}}{7m}+\cfrac{5q\bigl[\vec{v}_C\times\vec{B}\bigr]}{7m}-\cfrac{q\vec{B}\bigl({~}\vec{r}\cdot\vec{\omega}{~}\bigr)}{7m} $$Note, that still force lie in plane $xz$, which means that the motion is one-dimensional, since initially the ball was motionless

\end{solcard}

\begin{solcard}{B6}{Find the projection of the acceleration of the center $a_{Cx}$ of the ball onto the axis $x$ depending on the position of its center $x_C$.}{0.50}

Projecting the acceleration of the ball center onto the $x$ axis: $$a_{Cx}=\cfrac{5g\sin\alpha}{7}+\cfrac{qBr\omega_z}{7m} $$Substitute expression for $\omega_z$, we obtain:

\par\smallskip\noindent\textbf{Answer:} $$a_{Cx}=\cfrac{5g\sin\alpha}{7}-\cfrac{q^2B^2x}{14m^2} $$
\end{solcard}

\begin{solcard}{B7}{Determine the dependence of the coordinate $x_C$ of the center of the ball on time $t$. Express your answer in terms of $g$, $\alpha$, $q$, $m$, $B$ and $t$.}{0.70}

Introducing the variable $x'=x-10g\sin\alpha\left(\cfrac{m}{qB}\right)^2$, we get: $$\ddot{x}'=-\left(\cfrac{qB}{m}\right)^2\cfrac{x'}{14} $$This is the equation of harmonic oscillations with cyclic frequency $\omega_0=qB/m\sqrt{14}$. Then general solution for $x(t)$: $$x(t)=10g\sin\alpha\left(\cfrac{m}{qB}\right)^2+A\cos(\omega_0t+\varphi_0) $$Account for initial condition: $$ \begin{cases} x(0)=0\\ \dot{x}(0)=0 \end{cases} \Rightarrow \begin{cases} \varphi_0=\pi\\ A=10g\sin\alpha\left(\cfrac{m}{qB}\right)^2 \end{cases} $$For for $x(t)$ we have:

\par\smallskip\noindent\textbf{Answer:} $$x(t)=10g\sin\alpha\left(\cfrac{m}{qB}\right)^2\left(1-\cos\left(\cfrac{qBt}{m\sqrt{14}}\right)\right) $$
\end{solcard}

\begin{solcard}{B8}{Describe the evolution of the angular velocity of the ball $\vec{\omega}$. As an answer, indicate the hodograph of the angular velocity vector (in the coordinate system $xyz$), that is, a curve describing the position of the end of the vector with its fixed beginning. Indicate all characteristic values in the figure. Express them using $m$, $g$, $q$, $B$, $\alpha$ and $r$.}{0.50}

The angular velocity vector of the ball always lies in the $yz$ plane. For $\omega_z$ we have: $$\omega_z(t)=-\cfrac{qBx(t)}{2mr}=-\cfrac{5mg\sin\alpha}{qBr}\left(1-\cos\left(\cfrac{qBt}{m\sqrt{14}}\right)\right) $$For $\omega_y$ have: $$\omega_y(t)=\cfrac{v_{Cx}(t)}{r}=\cfrac{5mg\sin\alpha}{qBr}\sqrt{\cfrac{2}{7}}\sin\left(\cfrac{qBt}{m\sqrt{14}}\right) $$Hence: $$\sin^2\left(\cfrac{qBt}{m\sqrt{14}}\right)+\cos^2\left(\cfrac{qBt}{m\sqrt{14}}\right)=\left(\cfrac{\omega_y}{\cfrac{5mg\sin\alpha}{qBr}\sqrt{\cfrac{2}{7}}}\right)^2+\left(1+\cfrac{\omega_z}{\cfrac{5mg\sin\alpha}{qBr}}\right)^2=1 $$From the obtained relation it can be seen that the hodograph of the angular velocity vector is an ellipse with semi-axis $\omega_{0z}$ and $\omega_{0y}$, equal: $$\omega_{0z}=\cfrac{5mg\sin\alpha}{qBr}\qquad \omega_{0y}=\cfrac{5mg\sin\alpha}{qBr}\sqrt{\cfrac{2}{7}} $$

\begin{figure}[H]
  \centering
  \includegraphics[width=0.72\textwidth]{T3_S_EN_assets/fig1.jpg}
\end{figure}
\end{solcard}

\begin{solcard}{B9}{Find the maximum speed $v_{max}$ among all the elements of the ball at the moment when the center speed reaches its maximum. Express your answer in terms of $m$, $g$, $q$, $B$ and $\alpha$.}{0.60}

Let $\omega$ be the magnitude of the angular velocity forming an angle $\varphi$ with the plane.  The maximum speed is the point whose distance to the instantaneous axis of rotation of the ball is maximum. Further, it makes sense to consider only the section lying in the $yz$ plane. The figure shows that the maximum speed of the ball element is achieved at a perpendicular to the instantaneous axis of rotation passing through the center and is equal to: $$v_{max}=\omega r(1+\cos\varphi) $$


For the angular velocity components we have: $$\omega_y=\omega\cos\varphi\qquad \omega_z=-\omega\sin\varphi $$From the equation ellipse obtain: $$\cfrac{\omega^2\cos^2\varphi}{\omega^2_{0y}}+\cfrac{(\omega_{0z}-\omega\sin\varphi)^2}{\omega^2_{0z}}=1\Rightarrow{\omega=\cfrac{2\omega_{0z}\sin\varphi}{\sin^2\varphi+\cfrac{\omega^2_{0z}\cos^2\varphi}{\omega^2_{0y}}}} $$Hence for maximum velocity have: $$v_{max}=\cfrac{2\omega_{0z}r\sin\varphi(1+\cos\varphi)}{\sin^2\varphi+\cfrac{\omega^2_{0z}\cos^2\varphi}{\omega^2_{0y}}}=\cfrac{2\omega_{0z}r\tan\varphi(1+\sqrt{1+\tan^2\varphi})}{\tan^2\varphi+\cfrac{\omega^2_{0z}}{\omega^2_{0y}}} $$in consider moment $\tan\varphi=\omega_{0z}/\omega_{0y}$, therefore: $$v_{max}=\omega_{0y}r\left(1+\sqrt{1+\cfrac{\omega^2_{0z}}{\omega^2_{0y}}}\right) $$Substitute $\omega_{0y}$ and $\omega_{0z}$, we find:

\par\smallskip\noindent\textbf{Answer:} $$v_{max}=\cfrac{5mg\sin\alpha(3+\sqrt{2})}{qB\sqrt{7}} $$
\begin{figure}[H]
  \centering
  \includegraphics[width=0.72\textwidth]{T3_S_EN_assets/fig2.jpg}
\end{figure}
\end{solcard}

\begin{solcard}{B10}{At what minimum friction coefficient $\mu_{min}$ of a ball on a plane is motion possible without slipping? Express your answer in terms of $g$, $\alpha$, $q$, $m$ and $B$.}{0.60}

We have shown that the motion is one-dimensional, which means that the Lorentz force is zero at any moment. Then the normal reaction force of the plane $N$ at any moment is equal to $$N=mg\cos\alpha $$We determine the friction force $F_x$ from the theorem on motion of the center of mass: $$F_x=ma_{Cx}-mg\sin\alpha=mg\sin\alpha\left(\cfrac{5}{7}\cos\left(\cfrac{qBt}{m\sqrt{14}}\right)-1\right)\geq{-\cfrac{12mg\sin\alpha}{7}} $$from which friction force $F_x$ always be located in following range: $$-\cfrac{2mg\sin\alpha}{7}\geq{F_x}\geq{-\cfrac{12mg\sin\alpha}{7}} $$The minimum possible coefficient of friction is determined expression: $$\mu_{min}=\cfrac{|F_x|_{max}}{N} $$from:

\par\smallskip\noindent\textbf{Answer:} $$\mu_{min}=\cfrac{12\tan\alpha}{7} $$
\end{solcard}

\begin{solcard}{C1}{Determine the distance $S_{i{,}i+1}$ between the positions of the center of the ball at the moments of the $i$th and $(i+1)$th stops. Express your answer using $i$, $g$, $k$, $\alpha$, $q$, $m$ and $B$.}{0.80}

Since during rolling there is only a redistribution of the normal reaction force, the equation of the dynamics of rotational motion relative to the center of mass remains true, and the theorem on the movement of the center of mass must be changed:


The phase diagram of each half of the oscillations is an ellipse, shifted from position $x_0$ by the amount $\pm{\Delta{x}}$, where: $$x_0=10g\sin\alpha\left(\cfrac{m}{qB}\right)^2 $$$$\Delta{x}=10g\sin\alpha\left(\cfrac{m}{qB}\right)^2\cfrac{k}{r\tan\alpha} $$


From the phase diagram it is clear that if $a_i$ is the amplitude of half of the oscillation with number $i$, then the amplitude of half of oscillation $a_{i+1}$ with number $i+1$ is determined by the expression: $$a_{i+1}=a_i-2\Delta{x} $$Account for expression for $a_1$: $$a_1=x_0-\Delta{x} $$have: $$a_{i+1}=x_0-(1+2i)\Delta{x} $$Then for $S_{i{,}i+1}$ have: $$S_{i{,}i+1}=2a_{i+1} $$or:

\par\smallskip\noindent\textbf{Answer:} $$S_{i{,}i+1}=2\cdot 10g\sin\alpha\left(\cfrac{m}{qB}\right)^2\left(1-\cfrac{k(1+2i)}{r\tan\alpha}\right) $$
\begin{figure}[H]
  \centering
  \includegraphics[width=0.72\textwidth]{T3_S_EN_assets/fig3.jpg}
\end{figure}
\end{solcard}

\begin{solcard}{C2}{Determine the range of possible positions $x_{C\text{to}}$ of the last stop of the ball's center. Express your answer in terms of $g$, $k$, $\alpha$, $q$, $m$ and $B$.}{0.20}

Rolling stops if at the stopping point $$|F_\text{roll}|\leq kmg\cos\alpha/r $$This correspond following coordinate $x$:

\par\smallskip\noindent\textbf{Answer:} $$x\in\biggl[10g\sin\alpha\left(\cfrac{m}{qB}\right)^2\left(1-\cfrac{k}{r\tan\alpha}\right){;}10g\sin\alpha\left(\cfrac{m}{qB}\right)^2\left(1+\cfrac{k}{r\tan\alpha}\right)\biggr] $$
\end{solcard}

\begin{solcard}{C3}{Find the time $\tau$ of the ball's movement. Express your answer in terms of $g$, $k$, $\alpha$, $q$, $m$ and $B$.}{0.50}

Oscillations occur until the next stopping point is in the range $x_\text{stop}\in\bigl[-\Delta{x}{;}\Delta{x}\bigr]$, which leads to the expression for the number of $N$ oscillations: $$N=\Biggl[\cfrac{a_1}{2\Delta{x}}\Biggr]+1 $$Since $\tan\alpha \gg k/r$, to good accuracy we have: $$N=\cfrac{r\tan\alpha}{2k} $$For time motion have: $$t=\cfrac{\pi N}{\omega_0} $$or:

\par\smallskip\noindent\textbf{Answer:} $$t=\cfrac{\pi mr\tan\alpha}{kqB}\sqrt{\cfrac{7}{2}} $$
\end{solcard}

\begin{solcard}{C4}{What amount of heat $Q$ was released during the movement of the ball? Express your answer in terms of $g$, $k$, $\alpha$, $q$, $m$ and $B$.}{1.00}

First solution: Let's use the law of conservation of energy: $$Q=-\Delta{W}_p-E_k $$ where $\Delta{W}_p$ - change potential energy system, and $E_k$ is its kinetic energy. For $\Delta{W}_p$ have: $$\Delta{W}_p=-mg\sin\alpha x_0=-10mg^2\sin^2\alpha\left(\cfrac{m}{qB}\right)^2 $$ For finding kinetic energy determinable final angular velocity spinning sphere: $$E_k=\cfrac{I\omega^2_z}{2}=\cfrac{Iq^2B^2x^2_0}{8m^2r^2}=\cfrac{q^2B^2}{20m}\cdot 100g^2\sin^2\alpha\left(\cfrac{m}{qB}\right)^4=5mg^2\sin^2\alpha\left(\cfrac{m}{qB}\right)^2 $$


$$Q=-\Delta{W}_p-E_k $$


where $\Delta{W}_p$ is the change in the potential energy of the system, and $E_k$ is its kinetic energy.


For $\Delta{W}_p$ we have:


$$\Delta{W}_p=-mg\sin\alpha x_0=-10mg^2\sin^2\alpha\left(\cfrac{m}{qB}\right)^2 $$


To find the kinetic energy, we determine the final angular velocity of the ball's rotation:


$$E_k=\cfrac{I\omega^2_z}{2}=\cfrac{Iq^2B^2x^2_0}{8m^2r^2}=\cfrac{q^2B^2}{20m}\cdot 100g^2\sin^2\alpha\left(\cfrac{m}{qB}\right)^4=5mg^2\sin^2\alpha\left(\cfrac{m}{qB}\right)^2 $$


First solution: explicitly calculate the work done by the rolling friction force. For it we have: $$A_i=-\cfrac{kN}{r}\cdot{2A_i}=-\cfrac{2kmg\cos\alpha}{r}\cdot{a_i} $$ After substitution $a_i$ obtain expression for work: $$A=-\cfrac{20mg^2\sin\alpha\cos\alpha k}{r}\left(\cfrac{m}{qB}\right)^2\sum\limits_{i=1}^N \left(1+\cfrac{k}{r\tan\alpha}-\cfrac{2ik}{r\tan\alpha}\right)$$ Sum, obtain: $$A=-\cfrac{20mg^2\sin\alpha\cos\alpha k}{r}\left(\cfrac{m}{qB}\right)^2\left(\left(1+\cfrac{k}{r\tan\alpha}\right)N-\cfrac{N(N+1)k}{r\tan\alpha}\right)$$ Substitute $N$ and account for smallness, obtain: $$A=-\cfrac{20mg^2\sin\alpha\cos\alpha k}{r}\left(\cfrac{m}{qB}\right)^2\left(\cfrac{r\tan\alpha}{2k}-\cfrac{r\tan\alpha}{4k}\right)=-5mg^2\sin^2\alpha\left(\cfrac{m}{qB}\right)^2 $$ Amount heat $Q=-A$, from which:


$$A_i=-\cfrac{kN}{r}\cdot{2A_i}=-\cfrac{2kmg\cos\alpha}{r}\cdot{a_i} $$


After substituting $a_i$ we get the expression for work:


$$A=-\cfrac{20mg^2\sin\alpha\cos\alpha k}{r}\left(\cfrac{m}{qB}\right)^2\sum\limits_{i=1}^N \left(1+\cfrac{k}{r\tan\alpha}-\cfrac{2ik}{r\tan\alpha}\right)$$


Summing up, we get:


$$A=-\cfrac{20mg^2\sin\alpha\cos\alpha k}{r}\left(\cfrac{m}{qB}\right)^2\left(\left(1+\cfrac{k}{r\tan\alpha}\right)N-\cfrac{N(N+1)k}{r\tan\alpha}\right)$$


Substituting $N$ and taking into account the smallness, we get:


$$A=-\cfrac{20mg^2\sin\alpha\cos\alpha k}{r}\left(\cfrac{m}{qB}\right)^2\left(\cfrac{r\tan\alpha}{2k}-\cfrac{r\tan\alpha}{4k}\right)=-5mg^2\sin^2\alpha\left(\cfrac{m}{qB}\right)^2 $$


The amount of heat $Q=-A$, from where:

\par\smallskip\noindent\textbf{Answer:} $$Q=5mg^2\sin^2\alpha\left(\cfrac{m}{qB}\right)^2 $$
\end{solcard}

\begin{solcard}{D1}{Determine the modulus of the torque of the rotational friction force $M'$. Express your answer in terms of $\mu$, $m$, $g$, $\alpha$ and $r_0$.}{0.40}

The reaction force $N$ is uniformly distributed over the area of a circle with radius $r_0$, therefore: $$\cfrac{dN}{dS}=\cfrac{mg\cos\alpha}{\pi{r}^2_0} $$Then for moment force friction have: $$dM'=\cfrac{\mu dN}{dS}\cdot{2\pi{r}dr}\cdot{r}\Rightarrow{M'=\cfrac{\mu dN}{dS}\cdot{\cfrac{2\pi{r}^3_0}{3}}} $$from:

\par\smallskip\noindent\textbf{Answer:} $$M'=\cfrac{2\mu mgr_0\cos\alpha}{3} $$
\end{solcard}

\begin{solcard}{D2}{At what point in time $t_1$ will the ball start spinning? Express your answer in terms of $\mu$, $r_0$, $m$, $q$, $B$, $r$ and $\alpha$.}{0.50}

Since the arising torque of spinning friction is directed along the $z$ axis, the expression for the friction force $\vec{F}$, and, as a consequence, the theorem on the movement of the center of mass remains true: $$a_{Cx}=\cfrac{5g\sin\alpha}{7}+\cfrac{qBr\omega_z}{7m} $$in equation dynamics rotational motion necessarily introduce change: $$I\dot{\omega}_z=M_\text{fr}-\cfrac{qr^2B\omega_y}{5} $$where $M_\text{tr}$ - arise torque friction spinning. Until the velocity of the center of the spheree not reach value $$v_0=\cfrac{2M'mr}{IqB}=\cfrac{10\mu mgr_0\cos\alpha}{3qBr} $$spthe sphere moves with uniform acceleration and does not spin. Then expression for time $t_1$ following: $$t_1=\cfrac{7v_0}{5g\sin\alpha} $$or:

\par\smallskip\noindent\textbf{Answer:} $$t_1=\cfrac{14\mu mr_0}{3qBr\tan\alpha} $$
\end{solcard}

\begin{solcard}{D3}{Construct a high-quality graph $\omega_z(t)$. On the graph, indicate all the characteristic values, features and asymptotes. Express the characteristic values in terms of $m$, $g$, $\alpha$, $q$, $B$ and $r$.}{1.50}

To analyze the subsequent movement, we differentiate the expression for acceleration $a_{Cx}$: $$\dot{a}_{Cx}=\ddot{v}_{Cx}=\cfrac{qBr\dot{\omega}_z}{7m}=\cfrac{qBr}{7m}\left(\cfrac{M'}{I}-\cfrac{qBv_x}{2mr}\right) $$Introducing quantity $v_1=v_{Cx}-v_0$, obtain: $$\ddot{v}_1=-\left(\cfrac{qB}{m}\right)^2\cfrac{v_1}{14}\Rightarrow{v_1(t_1)=A\sin(\omega_0\tau+\varphi_0)} $$where time $\tau=t-t_1$ - time, elapsed from moment beginning/origin spinning. Use initial condition: $$ \begin{cases} v_1(0)=A\sin\varphi_0=0\\ \dot{v}_1(0)=\omega_0A\cos\varphi_0=\cfrac{5g\sin\alpha}{7} \end{cases} \Rightarrow v_{Cx}(\tau)=v_0+\cfrac{5g\sin\alpha\sin\omega_0\tau}{7\omega_0} $$Substituting into the equation dynamics rotational motion relatively axis $z$: $$\dot{\omega}_z=-\cfrac{qB}{2mr}\cdot{\cfrac{5g\sin\alpha\sin\omega_0\tau}{7\omega_0}}\Rightarrow{\omega_z(\tau)=-\cfrac{qB}{2mr}\cdot{\cfrac{5g\sin\alpha(1-\cos\omega_0\tau)}{7\omega^2_0}}} $$in the original terms: $$\omega_z(\tau)=-\cfrac{5mg\sin\alpha}{qBr}\left(1-\cos\left(\cfrac{qB\tau}{m\sqrt{14}}\right)\right) $$From the resulting expression it is clear that at the moment the spinning stops, the ball repeats it again according to the same law, since both the speed of the center and the angular speed of the spinning change with the same period. This can be interpreted as continuous spinning. Then the graph looks like the one shown in the figure below:

\begin{figure}[H]
  \centering
  \includegraphics[width=0.72\textwidth]{T3_S_EN_assets/fig4.jpg}
\end{figure}
\end{solcard}

\begin{solcard}{D4}{Construct a high-quality graph $x_C(t)$. On the graph, indicate all the characteristic values, features and asymptotes. Characteristic values can be expressed through the variables you enter.}{0.60}

We have already mentioned that before the start of slipping, during time $t_1$ the ball moves uniformly accelerated and manages to cover the distance: $$x_1=\cfrac{7v^2_0}{10g\sin\alpha}=\cfrac{70\mu^2m^2gr^2_0\cos^2\alpha}{9q^2B^2r^2\sin\alpha} $$We determine dependence $x(\tau)$: $$x(\tau)=x_1+\int\limits_0^{\tau}\left(v_0+\cfrac{5g\sin\alpha\sin\omega_0\tau}{7\omega_0}\right)d\tau=x_1+v_0\tau+\cfrac{5g\sin\alpha(1-\cos\omega_0\tau)}{7\omega^2_0} $$ Let's depict this on the graph:

\begin{figure}[H]
  \centering
  \includegraphics[width=0.72\textwidth]{T3_S_EN_assets/fig5.jpg}
\end{figure}
\end{solcard}

\end{document}